\section{Response, effect size, elasticity and interaction}
Four ideas recur in the figures of this chapter. The \kw{response} of an output $y$ to an input $x$ is the curve $y(x)$ with other factors averaged. The \kw{effect size} of a factor is the change of the output between its low and high levels as a share of the mean, $\mathrm{effect}=100\,(\bar y_{hi}-\bar y_{lo})/|\bar y|$; it allows factors of different units to be put on one scale. The \kw{elasticity} $\varepsilon=\mathrm d\ln y/\mathrm d\ln x$ is the per cent change of the output for a one per cent change of the input; $\varepsilon=1$ is proportional, $\varepsilon=0.3$ weak and $\varepsilon>1$ amplifying. An \kw{interaction} exists when the effect of one factor depends on the level of another; it is seen as non-parallel lines in an interaction plot or as curved contours in a response surface.

\section{Non-linearity, thresholds and regimes}
The soil--atmosphere system is not linear: a change of the input does not give a proportional change of the output. A \kw{threshold} is a value of the input at which the response changes in slope; it is found by fitting two straight lines that join at a break point $x_b$ and choosing $x_b$ that gives the smallest sum of squares. A \kw{regime} is a part of the factor space in which one process dominates: below a rainfall the evaporation dominates the loss, above it the drainage does; the boundary between the two regimes is where the two are equal, and it is the line drawn on the regime map. Near a threshold, a small change of the input can change the strategy: this is the reason why the averages over all conditions are poor guides to a field in a particular condition.

\section{Trade-offs and the Pareto front}
A management has several aims that may conflict: a high yield, a high irrigation efficiency, a low drainage, a low nitrate loss. A choice $r$ \kw{dominates} a choice $r'$ if it is at least as good in every aim and strictly better in one. The set of choices that no other choice dominates is the \kw{Pareto front}; every point on it is a trade-off in which an improvement in one aim can come only by a loss in another. The choice among the points of the front is not a matter of calculation but of the values of the manager and of the regulations (a limit on nitrate in the drainage, for example).

\section{Controllability}
The key figure introduces the \kw{controllability} of an outcome by a lever: the range of the outcome (its 10th to 90th percentile) over the values of the lever, as a share of the median. A high value means that the manager has influence; a low value means that the soil or the climate decides. The finding of the study is that the controllability is set by the soil pores: in a soil that drains freely the water goes down whatever the manager does, and in a soil that holds the water, the levers (trigger, dose, mulch) work. The rank correlation between the controllability and the available water of the ten soils measures how strong the link is.

\section{Nitrate concentration and the limit for drinking water}
The concentration of nitrogen in the water that drains is the nitrogen lost divided by the water lost: $c=N_{leach}/D\times100$ in mg\,N\,L$^{-1}$ when $N_{leach}$ is in kg\,N\,ha$^{-1}$ and $D$ in mm. The limit for drinking water in the European Union and in the guidelines of the World Health Organization is 50 mg of nitrate per litre, which is 11.3 mg of nitrogen. A given loss of nitrogen gives a lower concentration if the drainage is larger: the dilution may hide a problem, because it is the total load that goes to the sea. Both the concentration and the load are therefore given.
